% Section: Conclusion
\section{Conclusion}
\label{sec:conclusion}

Deployed agentic systems often apply a single fixed strategy to every query, even though
different queries demand different reasoning procedures and incur different execution costs.
We study \textbf{pre-inference agent-strategy routing} under a performance--cost objective:
given a query, select one strategy from a fixed pool before any strategy runs.

\textbf{Adaptive Agent Routing (AAR)} addresses this problem by combining planner-derived
procedural complexity with semantic query embeddings and training a router with
\textbf{cost-soft supervision} that preserves mass over cheaper correct strategies rather
than collapsing to a single hard winner.

On a held-out routing benchmark, \textbf{cost-soft supervision} contributes the largest
reduction in utility regret (agent regret $0.361\to 0.192$).
Planner-derived complexity further improves routing when \textbf{combined} with semantic
features; complexity-only representations underperform embedding-only baselines, while the
fused configuration attains the lowest total regret ($0.198$).
Under executed benchmark transfer, results are competitive but mixed and do not uniformly
outperform heuristic routing across families.

Future work includes improving cross-workload transfer, distilling planner-derived complexity
into cheaper features, expanding the strategy and backbone pools, and explicit
characterization of the accuracy--cost Pareto frontier beyond scalar utility with fixed~$\lambda$.
